\chapter{ConstitutiveIntegrator as the Algorithmic Injection Point}

\section{Motivation}

After introducing \texttt{ConstitutiveState}, the constitutive subsystem still
had one important ambiguity: the algorithmic step used to evaluate and commit a
constitutive site was semantically present, but it was still described under
the generic name \texttt{IntegrationStrategy}.  That name was operationally
useful, but it was too broad.  The relevant object at this level is not any
numerical integration procedure; it is the \emph{local constitutive
integrator}.

This chapter makes that distinction explicit by introducing
\texttt{ConstitutiveIntegrator}.

\section{Local Architectural Split}

The constitutive subsystem is now described by four roles:

\begin{enumerate}
  \item \texttt{ConstitutiveRelation}: the physical local law;
  \item \texttt{ConstitutiveState}: the state attached to the site;
  \item \texttt{ConstitutiveIntegrator}: the algorithm that queries and commits
        the site;
  \item \texttt{Material<Policy>}: the erased runtime handle that injects the
        chosen integrator into heterogeneous containers.
\end{enumerate}

This is not yet the final fully factored architecture, because many current
nonlinear laws still keep irreversible variables internally.  However, the
public and internal semantics now reflect the intended long-term split:

\begin{center}
\texttt{ConstitutiveLaw + ConstitutiveState + ConstitutiveIntegrator}
\end{center}

\section{Why the Injection Point Stays in \texttt{Material<Policy>}}

The local algorithm is injected at the \texttt{Material<Policy>} boundary on
purpose.

\begin{itemize}
  \item It preserves the existing type-erased heterogeneous storage model.
  \item It avoids contaminating each constitutive relation with multiple
        algorithmic variants prematurely.
  \item It keeps the hot path statically typed inside each homogeneous material
        model, while the runtime polymorphism remains confined to the already
        existing erased handle.
  \item It naturally matches the user-facing syntax:
\begin{lstlisting}[language=C++]
Material<ThreeDimensionalMaterial>{
    J2PlasticMaterial3D{E, nu, sigma_y, H},
    EmbeddedInelasticConstitutiveIntegrator{}
};
\end{lstlisting}
  \item It opens the door to future variants such as substepping, line-search
        local updates, operator splits, or explicit/implicit constitutive
        drivers without changing element kernels.
\end{itemize}

\section{Current Concrete Integrators}

Two semantic integrators are now made explicit:

\begin{itemize}
  \item \texttt{PassThroughIntegrator}: path-independent evaluation.  It reads
        stress and tangent directly from the constitutive relation and keeps
        \texttt{commit()} as a no-op for backward-compatible elastic behaviour.
  \item \texttt{EmbeddedRelationIntegrator}: compatibility bridge for current
        nonlinear laws.  It evaluates response and tangent from the relation,
        and during \texttt{commit()} it advances the relation through its own
        \texttt{update()} method while also updating the external
        \texttt{ConstitutiveState}.
\end{itemize}

The legacy names remain as aliases:

\begin{itemize}
  \item \texttt{ElasticUpdate} $\equiv$ \texttt{PassThroughIntegrator}
  \item \texttt{InelasticUpdate} $\equiv$ \texttt{EmbeddedRelationIntegrator}
\end{itemize}

This keeps the codebase source-compatible while allowing the semantic cleanup
to proceed.

\section{Why This Matters}

The key architectural benefit is that the current code can now express the
following distinction:

\begin{itemize}
  \item the constitutive \emph{law} is responsible for the physics;
  \item the constitutive \emph{integrator} is responsible for the local
        algorithmic realization of that physics inside a time/load step.
\end{itemize}

That distinction is essential for future nonlinear concrete and fracture work.
For example, the same inelastic law may later be integrated by:

\begin{itemize}
  \item an embedded monolithic return-mapping;
  \item a generic backward-Euler driver;
  \item a substepping wrapper;
  \item a viscoplastic regularized driver;
  \item a dynamic operator split coupled to finite-kinematics update rules.
\end{itemize}

\section{Custom Integrators}

The refactor also validates a more important property than the naming itself:
\texttt{Material<Policy>} is now an actual injection point for custom local
algorithms.

The new regression test introduces a custom integrator that intentionally
updates only the external constitutive-state carrier but does not evolve the
relation's internal irreversible variables.  The result is a predictable change
in post-commit behaviour, which proves that the algorithmic policy is not
hard-wired to the constitutive relation.

This is exactly the kind of freedom required for future scale-reconstruction
pipelines and multi-algorithm benchmarking.

\section{What Has Not Been Completed Yet}

This step is a semantic and architectural split, not a full behavioural one.
The following limitations remain:

\begin{itemize}
  \item many inelastic relations still own their irreversible state internally;
  \item the integrator does not yet receive an explicit algorithmic state
        object distinct from the constitutive relation;
  \item the current bridge integrator still relies on
        \texttt{relation.update(k)} for nonlinear constitutive evolution.
\end{itemize}

Therefore, the next meaningful refactor should move the local evolution data
out of the constitutive relation and into explicit state objects wherever that
is mathematically coherent.

\section{Immediate Architectural Interpretation}

After this step, the recommended reading of the subsystem is:

\begin{itemize}
  \item \texttt{ConstitutiveRelation}: local physical constitutive model;
  \item \texttt{ConstitutiveState}: state attached to the site;
  \item \texttt{ConstitutiveIntegrator}: local algorithmic driver;
  \item \texttt{MaterialInstance}: typed constitutive site;
  \item \texttt{Material<Policy>}: erased constitutive handle with injected
        local integrator.
\end{itemize}

This is a materially more precise ontology than the previous generic
``strategy'' language, and it is better aligned with the future goals of
inelastic reinforced concrete, dynamics, large displacements, and
continuum-scale damage or fracture.
